% abstract.tex
This is Paper~II of the A=1 Discrete Causal Lattice series. The
central thesis is to prove (or characterise the obstruction to)
Eq.~(137) of Paper~I: that the automorphism group of the bipartite
octahedral causal lattice $\mathcal{T}_\diamond^3$ under the unity
constraint $\mathcal{A}=1$ is the Standard Model gauge group plus
Lorentz, $SO(3,1) \times SU(3) \times SU(2) \times U(1)$, with real
dimension $18 = 6+8+3+1$.  We establish on the existing per-site
$\mathbb{C}^2 = (\psi_R, \psi_L)$ amplitude that $SO(3,1) \times U(1)$
acts as a dim-7 sub-product (verified symbolically), that the proposed
per-site $SU(2)$ generators on this carrier are not independent of
the Lorentz rotations, and that the RGB sublattice contributes only
$\mathbb{Z}_3 \subset SU(3)$.  We extend the framework to the per-site
amplitude $\mathbb{C}^{12} = \mathbb{C}^2 \otimes \mathbb{C}^2 \otimes
\mathbb{C}^3$ (chirality $\otimes$ weak isospin $\otimes$ colour),
verify that the four conjecture factors commute pairwise on this
carrier (dim 18), verify that a trivial tensor extension of the
bipartite tick rule preserves $\mathcal{A}=1$ and the global
$SU(2)_W \times SU(3)$ symmetries, and verify that the bipartite
Wilson plaquette $V_1, -V_2, -V_1, V_2$ is gauge-invariant.  The
remaining open subproblems --- exact equality versus containment in
$\mathrm{Aut}_\text{ext}$, SM-chirality coupling alignment, and the
$1/g^2$ prefactor for the non-abelian Wilson actions --- are stated
as the substance of the paper.  The paper closes with either a
derivation of the SM gauge group from a single conservation axiom on
a discrete substrate, or a precise statement of the obstruction; both
outcomes are concrete results.

\placeholder{(Trim and tighten as the paper develops; the audit table
in Section~\ref{sec:introduction} carries the canonical PASS/STUB
status for each sub-claim.)}
